\documentclass[12pt,a4paper,oneside]{article}
\usepackage{graphicx}
\usepackage{float} %for [H]


\usepackage{titlepic}
\usepackage[utf8]{inputenc}
\usepackage[left=1.1in,right=1.1in, top=1in, bottom= 1in]{geometry}
\usepackage[none]{hyphenat} % Avoids to go out of margin
\usepackage{subfiles}

% Font size of figure smaller than normal size:
\usepackage{caption}
\captionsetup[figure]{font=small}
\linespread{1.2}

% --------------------------------------------- %

\title{Homework 2}	                                    % Title
\author{Flavio Maiorana 2051396}				        % Authors separated by \\
\date{\today}								            % Date

\makeatletter
\let\thetitle\@title
\let\theauthor\@author
\let\thedate\@date
\makeatother

\begin{document}

\begin{titlepage}
	\centering
    \vspace*{0.5 cm}
    \includegraphics[scale = 0.75]{figures/SapienzaLogo.pdf}\\[1.0 cm]	% University Logo

    % \vspace*{-0.4cm}
    % \textsc{\large DIAG}\\[2.0 cm]	% Department Name
    % \vspace*{1cm}

    { \fontsize{20.74pt}{18.5pt}\selectfont\bfseries \thetitle \par } % Title

    \vspace*{0.25cm}
    \textsc{\Large Machine Learning}\\[0.5 cm] % Course Name

    \vspace*{2.6cm}
	% \begin{minipage}{0.4\textwidth} % 0.4
	% 	\begin{flushleft} \large
	% 		\textbf{Professors:}\\
	% 		Professor Name\\
	% 	\end{flushleft}
	% \end{minipage}~
	\begin{minipage}{0.3\textwidth} %0.4
		\begin{flushright} \large
		\begin{minipage}{1\textwidth}
		\begin{flushleft} \large
			\textbf{Students:} \\
			\theauthor
        \end{flushleft}
        \end{minipage}
		\end{flushright}
	\end{minipage}\\[3.85 cm]

    \vspace{2cm}
    \rule{\linewidth}{0.2 mm} \\[0.3 cm]
    \vspace*{-0.2cm}
    Academic Year 2023/2024
\end{titlepage}
\newpage


\section{Introduction}

First of all, it could be useful to gain some insight one how the dataaset is
made. The available data is already split into training and testing set. We
will consider the test split only in the evaluation phase (to prevent
data leakage).

\begin{verbatim}[Dataset for Training]
    N Examples: 6369
    N Classes: 5
    Classes: [0 1 2 3 4]
    - Class 0: 1000 (15.701051970482022)
    - Class 1: 1500 (23.551577955723033)
    - Class 2: 1500 (23.551577955723033)
    - Class 3: 2000 (31.402103940964043)
    - Class 4: 369 (5.7936881771078665)
\end{verbatim}

Some comments: the dataset is highly imbalanced. That needs to be addressed
accordingly. The main two problems to address are:

\begin{itemize}
    \item Model architecture
    \item Training process
    \begin{itemize}
        \item Class imbalance mitigation
        \item Optimization choice
        \item Performance metrics choice
    \end{itemize}
\end{itemize}

\section{Model Architecture}

The Model is a CNN, since we have to classify images. The first and foremost
choice regarding cnn model architecture is surely how deep it has to be and how
big the kernel need to be. Since the images are not that big, a reasonably small
cnn would be enough, since a too deep one would either overfit on training data,
or waste computational resources without gaining any improved performance.
Although, also a too small model would not be able to adequately catch the
features of the image. Also, kernels size is very important. A too big kernel
would compromise computational efficiency, but the size has to be adapted to the
tiype of features the convolution needs to to capture. To corroborate these
concepts we will compare some cnns with a different number of convolutional
layers. Each cnn has 2 fully connected layers at the end. Each list entry
indicates the dimension for the convolutional layer (from 1 to N). The kernel
size 0 indicates the number of channels of the input image. This training
results are all captured after 40 epochs.

\begin{itemize}
    \item N = 3
    \begin{verbatim}
        'channels': [3,8,16,32,64],
        'kernels': [5,5,5,5],
        'strides': [1,1,1,5],
        'pool_kernels': [2,2,2,2],
        'pool_strides': [2,2,2,2]

        precision    recall  f1-score   accuracy
          0.702     0.623     0.644      0.623 

        Total reward: 903.86
        Frames: 894
        Training time: 300.85
    \end{verbatim}
    \item N = 6 
    \begin{verbatim}
        'channels': [3,10,16,32,64,64,64,64,64],
        'kernels': [5,5,5,3,3,3,64,64],
        'strides': [1,1,1,1,1,1,1,1],
        'pool_kernels': [2,2,2,2,2,2],
        'pool_strides': [2,1,2,1,2,2],

        precision  recall  f1-score   accuracy
          0.699    0.638     0.653      0.638

        Tile points: 874.69
        Training time:
    \end{verbatim}
\end{itemize}

\subsection{Batch normalization layers}

\subsection{Dropout layer}

\subsection{Ensemble}

\newpage

\section{Training process}

There are some important issues to address, which
are to some extent independent from the model we use to classify. First of all,
to mitigate class imbalance we could apply:
\begin{itemize}
    \item Data augmentation (useful in any case to increase generalization)
    \item Different sampling techniques
    \begin{itemize}
        \item Upsampling/Downsampling
        \item Weighted resampling
    \end{itemize}
\end{itemize}

\subsection{Sampling techniques}

In order to mitigate class imbalance, one could try to rebalance the dataset by
using weighted resampling techniques.

\subsection{Data augmentation}

This technique is widely used with images. It exploits the invariance property
of images to enhance generalization properties of neural networks. In our case
it is particularly useful, because it allows us to "reuse" the same sample
multiple times and make it look new to the network by augmenting its features.
The transformation I used are these three from the torchvision package, plus a
random color jitter implemented separately. I also tried
RandomHorizontalFlipping, but I noticed actually worse conditions, maybe because
the images are sensitive to the horizontal orientation (for turning left and
right), so randomly flipping them would mean introducing noise to the labels.

\begin{verbatim}
    transforms.RandomInvert(),
    transforms.RandomErasing(),
    transforms.RandomRotation(degrees=5),
\end{verbatim}

\subsection{Weighted random resampling}

To rebalance classes, we could sample differently the training data, maybe
according to some weights. 

\section{Final Considerations}

In the end, after different trials I was able to achieve the completion of a
full lap through manual tuning of class weights 

\bibliographystyle{unsrt}
\bibliography{ref}
\nocite{*}  % to include references which were not cited

\end{document}